\documentclass[11pt]{article}
\usepackage[T1]{fontenc}
\usepackage[utf8]{inputenc}
\usepackage{lmodern,amsmath,amssymb,amsthm,mathtools,booktabs,array,longtable}
\usepackage[margin=1in]{geometry}
\usepackage{microtype,xurl}
\usepackage[colorlinks=true,linkcolor=blue,urlcolor=blue,citecolor=blue]{hyperref}
\hypersetup{pdftitle={Direct and centered Diophantine consistency equations},pdfauthor={dot (AI-assisted research for Vladimir Reshetnikov)},pdfsubject={Round twelve mathematical constructions and reproducible artifacts},pdfkeywords={Diophantine equations, ZFC consistency, Pell equations, coefficient height}}
\usepackage{enumitem}
\setlist{nosep}
\setlength{\emergencystretch}{3em}
\newcommand{\N}{\mathbb N}
\newcommand{\Z}{\mathbb Z}
\newcommand{\Con}{\operatorname{Con}}
\newtheorem{theorem}{Theorem}
\newtheorem{lemma}[theorem]{Lemma}
\title{Direct and centered Diophantine consistency equations}
\author{AI assisted research prepared by dot for Vladimir Reshetnikov}
\date{Round twelve\quad 24 October 2026}
\begin{document}
\maketitle
\begin{abstract}
We give a fully materialized ordinary quartic polynomial whose lack of
integer zeros is provably equivalent in ZFC to syntactic consistency of ZFC.
The construction starts directly from the fixed numerical C222 proof calculus,
uses a coefficient-height certificate for an unordered proof, and eliminates
one bounded family by explicit CRT, factorial, binomial, power and Pell
representations. It uses neither K215 machine simulation nor Jones's universal
polynomial. The result has 27,916 unrestricted integer variables, 758,275
nonzero monomials, and largest absolute coefficient 320,307,372,009,918.
A companion has 16,036 nonnegative variables. This is an architectural
completion, not a variable-count record.

Separately, integral square completion in a final radical norm gives four
10-integer cost--degree tradeoffs, down to degree 1,080,888. A cofinal
auxiliary Pell orbit yields an 11-positive-variable equation of degree 211
with 113 live arithmetic operations. Reversible norm-row reduction gives
21 unrestricted integers at degree 16. The complete literal equations,
actual fixed integer payloads, mathematical proof interfaces, independent
checkers and deterministic replay are supplied. None of these claims is a
new end-to-end Lean or Rocq certificate or a proof of consistency.
\end{abstract}
\tableofcontents
\newpage
\section{Results and the meaning of simplicity}
The target statement is the usual syntactic sentence
$\Con(\mathrm{ZFC})$: there is no finite ZFC proof of a fixed contradiction.
Throughout, $\N=\{0,1,2,\ldots\}$, whereas a positive variable ranges over
$\Z_{>0}$. For every completed equation identified below, ordinary mathematics
formalizable in ZFC proves
\[
 \Con(\mathrm{ZFC})\quad\longleftrightarrow\quad
 \neg\exists\boldsymbol x\in D^n\;P(\boldsymbol x)=0,
 \qquad D\in\{\N,\Z,\Z_{>0}\}.
\]
This biconditional does not establish either side. A fixed coefficient may
be enormous without being an unknown; a circuit register may abbreviate a
polynomial without being a quantified variable. The report keeps these
notions separate.

There is no single total order called ``simplest.'' The following completed
round-twelve points optimize different resources. A binary arithmetic operation
means one addition, subtraction or multiplication, including a square or a
scalar product; copying and reading a fixed integer leaf cost zero.
\begin{center}\small
\begin{tabular}{llrrl}
\toprule
Construction & Domain & Variables & Degree & Live operations\\
\midrule
Direct proof source & $\Z$ & 27,916 & 4 & expanded polynomial\\
Direct proof source & $\N$ & 16,036 & 4 & quadratic residual list\\
Centered auxiliary row & $\Z$ & 21 & 16 & not optimized\\
Two small Pell gaps & $\Z_{>0}$ & 11 & 211 & 113\\
Short guard and polarization & $\Z$ & 10 & 1,245,324 & 471\\
First radical center & $\Z$ & 10 & 1,179,788 & 485\\
Second radical center & $\Z$ & 10 & 1,114,252 & 496\\
Further integral correction & $\Z$ & 10 & 1,080,888 & 545\\
\bottomrule
\end{tabular}
\end{center}
For a compact evaluable ten-integer equation, the 471-operation point is the
lowest-cost delivered choice. If degree matters more, use one of the centered
alternatives. For a transparently completed direct route independent of the
machine/Jones architecture, use the large quartic. Neither recommendation
asserts a global optimum.

\subsection{What is new in this iteration}
The direct route was previously a proposal ending at an unmaterialized
compiler interface. Its final ordinary equation is now present, with every
coefficient and monomial, and each primitive interface independently audited.
The common witness bound, fixed binary height bound, affine-wire elimination,
and restriction of nonnegativity certificates to semantic inputs reduce the
actual compiled size. The mathematical ingredients of coefficient-height
encoding, CRT and Pell representation themselves are inherited.

In the small-equation branch, the new work is the source-aware short guard,
three explicit integral square completions, simultaneous transport of two
small Pell gaps with a cofinality proof, and the compatible auxiliary-norm row
operation for 21 integers at degree 16. A 12-positive-variable degree-70
100-operation circuit is also supplied, but it materializes a prior manuscript
formula and is not claimed as a new theorem.

\subsection{Comparison boundaries}
The repository comparison is pinned to
\href{https://github.com/VladimirReshetnikov/ProveIt/tree/9a33342a03d19cf5b061b2083c6b5d3f27f752ed}{ProveIt commit 9a33342a03d19cf5b061b2083c6b5d3f27f752ed}.
Six authenticated incoming archives were also inspected as inert data.
A later Library supplement, \emph{Pell Congruences}, is separate evidence,
not content attributed to that repository pin. It already contains
$(n,d)=(20,20),(19,28),(18,32),(17,88)$ over unrestricted integers.
Consequently our $(20,20)$ charts are alternatives, and our $(19,32)$ and
$(18,88)$ charts are dominated. Only the new $(21,16)$ point is promoted in
that low-degree comparison. The supplement's 9-positive-variable
1,158,228-degree 726-operation claim is distinct from the earlier
1,753-operation packet; it is not treated here as a new independently
reconstructed result. The repository's 9-positive-variable 440-operation
point has degree 1,845,940. Fewer variables, smaller degree, smaller circuits
and smaller numeral payloads remain distinct objectives.

The source audit records exact archive and retained-file hashes. Hashes
establish identity and provenance, not mathematical correctness. No upstream
program was executed during this research replay.


\section{A fully materialized equation from unordered proofs}
\label{sec:direct-materialized}

This section closes a different gap from the low-variable record search.
The height-certificate source previously ended at an unelaborated compiler
interface. We now give an ordinary polynomial with every coefficient and
monomial materialized. It is independent of the K215 transition table and
of Jones's universal polynomial. Its size is large: the purpose is a
complete, reproducible direct baseline, not a claim of a new variable record.

The recommended artifact is the literal ordinary polynomial $Q$ in
\path{direct/fixed_base/equation_integer.json}. It has 27,916 unrestricted
integer variables, exact degree four, 758,275 nonzero monomials, and
maximum absolute coefficient 320,307,372,009,918. In ZFC,
\[
 \operatorname{Con}(\mathrm{ZFC})\quad\longleftrightarrow\quad
 \neg\exists\boldsymbol{x}\in\mathbb Z^{27916}\;Q(\boldsymbol{x})=0.
\]
The natural-domain companion has 16,036 nonnegative variables. It is
specified exactly as the sum of the squared residuals in
\path{fixed_base/residuals_affine.json}; rebuilding also emits its
collected sparse polynomial \path{equation_affine.json}. Here
$\mathbb N=\{0,1,2,\ldots\}$. Each sparse entry
\texttt{[coefficient,[indices]]} denotes that coefficient times the product
of the listed variables; repeated indices denote repeated factors and
indices start at zero. All coefficients are literal integers, not bounded
search recipes. The compact archive retains the full integer polynomial
and exact residual/circuit data; rebuilding restores redundant expanded
intermediates and checks their recorded SHA256 hashes.

The derivation is mathematical; we do not claim a newly kernel-checked
end-to-end translation. We first explain the two-stream source and then
replace its height stream by a fixed binary bound. The final measurements
above include this reduction and affine-wire elimination.

\subsection{The fixed proof source}

The inherited file \path{fixed_base/horn_shared_source.json} specifies twelve
padded single-conclusion rules. Their local relation is
\[
 H(p,q,c;\boldsymbol w)=\prod_{j=0}^{11}
             \left(\sum_{r\in\mathcal R_j}r(p,q,c,\boldsymbol w)^2\right)=0.
\]
There are seventeen shared nonnegative witnesses. The residual lists
$\mathcal R_j$, including the fixed bundled axiom, are included literally.
Their expanded sums of squares are independently reproduced exactly by
the checker. The historical JSON field name
\path{quadratic_residuals} is not a degree assertion: the bundled branch
also contains quartic residuals. The total product has degree fifty. The compiled endpoint deliberately uses
this twenty-coordinate local source (three labels and seventeen Horn witnesses),
not the earlier twenty-three-coordinate degree-fourteen selector variant.
After adding six beta coordinates and three label-bound slacks, the bounded
row has twenty-nine witnesses. Reducing local polynomial degree need not
reduce the final quartic: each extra bounded witness adds arithmetic
certificates, whereas circuit flattening already makes the global degree four.
The first rule is truth $(0,0,0)$; zero premises pad axioms and one-premise
rules. Target code one is falsity.

For clarity, the calibration does not pass through execution of K215.
The repository theorem \path{c222_iff_zf_refutation} in
\path{KEndpoint.lean}, followed by
\path{zf_refutation_iff_zfc_refutation} in
\path{ConstructibleChoice.lean}, identifies derivability of target one
in this numerical calculus with a syntactic ZFC refutation. The exact
numeric constructors and twelve rules are in \path{KCode222.lean} and
\path{Code254.lean}. These are dependencies of the mathematical
calibration, not claims that the new Python compiler has been checked by Lean.

A finite family of valid rules has a \emph{flow certificate} for one when,
for every positive label $j$,
\[
 \#\{i:c_i=j\}=[j=1]+\#\{i:p_i=j\}+\#\{i:q_i=j\}.
 \tag{F}
\]
Derivability implies (F) by unfolding a proof tree. Conversely, if one is
not derivable, sum these balances over the nonderivable labels in the
finite support. Each valid single-conclusion rule contributes at most
zero: a nonderivable conclusion has a nonderivable premise. The right
side requires a positive net contribution at one, a contradiction.
This argument permits cycles and repeated premises but uses both
nonnegative multiplicities and the single-conclusion hypothesis.

\subsection{The coefficient-height certificate}

Given $N$ rule instances, put
\[
 A(X)=\prod_{i<N}(1+c_iX),\quad
 B(X)=(1+X)\prod_{i<N}(1+p_iX)(1+q_iX),\quad
 R_N=2\prod_{i<N}((1+p_i)(1+q_i)+c_i).
\]
Every coefficient of $A-B$ has absolute value at most $R_N$.
Consequently, for $M=R_N+1$, equality $A(M)=B(M)$ forces $A=B$:
the least nonzero coefficient of a nonzero polynomial vanishing at $M$
would be a nonzero multiple of $M$ of absolute value smaller than $M$.
The distinct roots $-1/j$ recover (F) from the polynomial identity.

This yields the two-accumulator source
\begin{align*}
&S_0=M+1,\quad S_N=1,\quad R_0=2,\quad M=R_N+1,\\
&S_{i+1}(1+c_iM)=S_i(1+p_iM)(1+q_iM),\\
&R_{i+1}=R_i((1+p_i)(1+q_i)+c_i),\qquad H(p_i,q_i,c_i;\boldsymbol w_i)=0.
\end{align*}
All factors are positive. Telescoping gives $A(M)=B(M)$; the second
accumulator supplies the coefficient bound. Conversely, expand a proof
tree in open-obligation order, compute the complete $R_N$ first, choose
$M=R_N+1$, and let $S_i$ be the product of the open-obligation factors.
There is no circular choice of the coefficient bound.

Encode the two sequences using a common beta scale:
\[
 \beta(x,y,i)=x\bmod(1+(i+1)y).
\]
Every pair of finite natural sequences admits such an encoding with
separate $x,z$ and one common $y$, by choosing $y$ large and divisible by
$N!$, then using the Chinese remainder theorem. The factorial divisibility
here is used only in the existence proof; it is not an additional
universally quantified condition in the final source.

The five external parameters are $N,x,z,y,M_0$, with $M=M_0+1$.
Four beta endpoint equations are imposed. In the one bounded row block,
the four beta values are replaced by four triples $(r,q,s)$ satisfying
\[
 x=qd+r,\qquad d=r+s+1,\qquad d=1+(i+1)y,
\]
with the appropriate code and index. Thus the row has exactly
$20+4\cdot3=32$ natural witnesses: twenty for labels and Horn witnesses,
and twelve for the beta graphs. There is a fixed signed polynomial
$P=e-f$ with $e,f$ having nonnegative coefficients such that the source is
\[
 \text{four endpoint equations}\quad\land\quad
 \forall i<N\;\exists\boldsymbol w\in\mathbb N^{32}\quad
 e(i,\boldsymbol w)=f(i,\boldsymbol w).
 \tag{S}
\]
For this two-stream comparison, $e,f$ can be represented as nonnegative
arithmetic DAGs. The delivered \path{fixed_base/row_source.json} instead
contains the twenty-nine-witness specialization described below. Both use
the following identities without expanding the large Horn product:
$(e-f)+(g-h)=(e+g)-(f+h)$ and
$(e-f)(g-h)=(eg+fh)-(eh+fg)$.

\subsection{One uniform witness bound and one CRT block}

This materialization uses a small improvement to the generic bounded
universal compiler: one common bound $U$ replaces thirty-two independent
coordinate bounds, and the factorial $(U+1)!$ is computed only once.
Set
\[
 T=(e(N,U,\ldots,U)+f(N,U,\ldots,U)+N)!,\qquad
 D_j=1+(j+1)T,\qquad B=\prod_{j<N}D_j.
\]
The bounded block in (S) is equivalent to the following finite
existential conditions, with encoded coordinates $a_1,\ldots,a_{32}$
and one encoded index $v$:
\begin{align*}
 B&\mid (U+1)!\binom{a_i}{U+1} &&(1\leq i\leq32),\\
 B&\mid 1+(v+1)T,\\
 e(v,\boldsymbol a)&\equiv f(v,\boldsymbol a)\pmod B.
\end{align*}

For soundness, fix $j<N$ and take a prime $p\mid D_j$. Since
$\gcd(p,T)=1$ and $T$ is the displayed factorial, $p$ exceeds its factorial
argument. Divisibility of each falling factorial supplies a residue
$w_i\leq U$ with $a_i\equiv w_i\pmod p$. This statement also handles
truncated natural subtraction: if a vanishing factor arises from
$a_i<t$, choose $w_i=a_i$. The index congruence gives $v\equiv j\pmod p$.
Hence $e(j,\boldsymbol w)\equiv f(j,\boldsymbol w)\pmod p$.
Monotonicity of nonnegative-coefficient polynomials bounds both values
strictly below $p$, so they are equal. For completeness choose $U$ above
all finitely many row witnesses. The $D_j$ are pairwise coprime because
$T$ is divisible by all positive differences below $N$; CRT supplies the
encoded columns and index. The argument is also valid for $N=0$.
The original source's endpoints themselves exclude $N=0$.

This is a uniform-bound specialization of the pinned theorem
\texttt{boundedForallCertificate}. It removes no semantic requirement and
needs no unproved new representation theorem.

\subsection{Eliminating every nonpolynomial operation}

The file \path{fixed_base/build_direct.py} implements the following explicit
identities. Remainders and quotients always have positive denominators
and are encoded by $n=qd+r$, $d=r+s+1$. Congruence modulo $d$, including
$d=0$, is encoded by $a+du=b+dv$ with two natural quotients. If
$Q=((N+1)T+1)^N+2$, then
\[
 B=\bigl(T^N N!\binom{Q+N}{N}\bigr)\bmod(TQ-1).
\]
The denominator is positive since $T\geq1$ and $Q\geq3$.
The binomial and factorial identities are
\begin{align*}
 \binom nk&=\left\lfloor\frac{(2^n+2)^n}{(2^n+1)^k}\right\rfloor
                \bmod(2^n+1),\\
 r&=(n+2)^{n+2},\quad C=\binom{r+n}{n},\qquad
 n!=\left\lfloor\frac{(r+1)^n+C-1}{C}\right\rfloor.
\end{align*}
They include all boundary cases, including $k>n$ and $n=0$.
The predecessor $C-1$ is represented by $C=C_-+1$; here $C>0$.
The pinned \path{FactorialRepresentation.lean} proves these identities,
and \path{BoundedForallRepresentation.lean} supplies the product formula.

Natural exponentiation is replaced by its fully explicit finite Pell
certificate, not by an oracle. Write
\[
 \mathcal N(a,x,y):\quad x^2+y^2-a^2y^2=1.
\]
The graph of $(x,y)=(X_n(a),Y_n(a))$, for $a>1$, is given by $n\leq y$
and the disjunction of $x=1,y=0$ with the existence of $u,w,s,t,b$ satisfying
\begin{gather*}
 \mathcal N(a,x,y),\quad \mathcal N(a,u,w),\quad \mathcal N(b,s,t),
 \quad b>1,\quad w>0,\\
 b\equiv1\pmod{4y},\quad b\equiv a\pmod u,\quad y^2\mid w,
 \quad s\equiv x\pmod u,\quad t\equiv n\pmod{4y}.
\end{gather*}
This is exactly the pinned finite Matiyasevich certificate in
\path{PellRepresentation.lean}. To represent $m=n^k$, use the disjunction
\[
 (k=0\land m=1)\ \lor\
 \left(k>0\land\left[(n=0\land m=0)\ \lor\
       (n>0\land\mathcal C(m,n,k))\right]\right),
\]
where $\mathcal C$ asserts existence of $w,a,t,z,x,y$ with the displayed
Pell-coordinate graph at $(a,k,x,y)$ and
\begin{gather*}
 x\equiv y(a\mathbin{\dot-}n)+m\pmod t,\qquad
 2an=t+n^2+1,\qquad m<t,\\
 n\leq w,\qquad k\leq w,\qquad \mathcal N(w+1,a,wz).
\end{gather*}
Here $\dot-$ is truncated natural subtraction. Its graph is encoded by
$(a=d+n)$ or $(a\leq n\land d=0)$. These formulas are the exact
\path{PowerRepresentation.lean} criterion. Strict and weak inequalities
use natural slack variables; divisibility uses a natural quotient.
Conjunctions are sums of nonnegative residual polynomials, and
disjunctions are their products. Consequently inactive branch witnesses
are harmless. In particular, no positivity assumption about an arbitrary
power base is hidden in the generated polynomial.

The delivered fixed-base call ledger contains 108 power certificates,
108 indexed Pell certificates, 33 binomial calls, three factorial calls,
and 72 division/remainder calls. The ledger records actual input and output
node indices and existential ranges for each call.

\subsection{Why the final polynomial is quartic}

All signed expressions are maintained as differences of natural
arithmetic circuits. Every addition or multiplication node receives a
fresh natural wire and a defining residual of degree at most two.
The final equality residuals are linear in the appropriate output wires.
The delivered polynomial is the sum of the squares of all these residuals.
A zero therefore enforces every circuit gate and every arithmetic
certificate. Conversely, any certificate uniquely supplies all gate
wires. Thus flattening preserves existence of zeros in both directions.
A nonzero squared multiplication residual ensures exact degree four;
the checker also verifies the degree of the fully collected polynomial.

Rebuilding emits \path{residuals.json} and \path{equation.json}, containing,
respectively, the complete residual list and its fully collected sum of
squares. Independent integer arithmetic reconstructs every gate and every
coefficient. The test suite also expands all twelve Horn branches exactly,
checks row evaluations against a separately written source expression,
checks all zero-base and zero-exponent boundaries, and includes a genuine
nonzero-exponent Pell witness for $1^1=1$. These finite checks support
artifact fidelity; the all-input equivalence rests on the proofs above.

\subsection{The fixed binary height bound}

The completed compiler makes a further source improvement measurable.
Choose an independent maximum label $m$, and impose
$p_i,q_i,c_i\leq m$ by three natural slack variables per row. Since
$m+1\leq2^m$ for every natural $m$,
\[
 A(1)\leq(m+1)^N\leq2^{mN},\qquad
 B(1)\leq2(m+1)^{2N}\leq2^{2mN+1}.
\]
It is therefore enough to set
\[
 M_0=2^{2mN+1},\qquad M=M_0+1.
 \tag{BH}
\]
The entire $R$ stream and its beta code disappear. There is now only the
$S$ stream, with endpoints $\beta(x,y,0)=M_0+2$ and
$\beta(x,y,N)=1$. Its local transition is unchanged. The row has
$20+6+3=29$ witnesses instead of thirty-two: twenty for Horn/labels,
six for two beta triples, and three label upper-bound slacks. One global
power certificate imposes (BH). Thus eight power certificates disappear
net, reducing 116 power calls to 108.

Completeness chooses a finite proof tree, then $N$ and a bound $m$ on
its labels, then $M_0$ by (BH), then the $S$ sequence, and finally its
beta encoding and the uniform CRT witness bound $U$. In particular,
$m$ is \emph{not} identified with $U$: the latter also bounds local
remainders and quotients depending on $M$. Such an identification would
introduce a circular, potentially impossible bound. Soundness uses the
same coefficient-height lemma and flow theorem as before. For $N=0$,
(BH) gives $M_0=2$ and the two endpoints are contradictory.

The folder \path{direct/fixed_base/} gives this source and its complete
polynomials. A second, slightly tighter but three-wire-larger version
uses $M_0=2(m+1)^{2N}$; it is an intermediate comparison, not needed
in the recommended release and not a distinct arithmetic theorem.

\subsection{Eliminating affine wires and allowing unrestricted integers}

Addition gates do not need existential output wires. Substitute their
affine expressions into the next multiplication residual. Scalar products
are affine too and are substituted likewise. Retain a wire only for a
multiplication whose expanded expression is genuinely quadratic in
previous retained wires and semantic inputs. Every residual still has
degree at most two. The substitution is triangular and uniquely
restores every omitted wire. All substituted affine forms have
nonnegative coefficients, so natural witnesses are preserved.

Applied to the fixed-binary-height source, this gives the recommended
natural polynomial (expanded by rebuilding as
\path{fixed_base/equation_affine.json}):
\begin{center}
\begin{tabular}{lr}
Natural variables & 16,036\\
Exact degree & 4\\
Monomials & 700,865\\
Semantic inputs & 3,960\\
Retained nonlinear wires & 12,076\\
Squared residuals & 12,366\\
Maximum absolute coefficient & 320,307,372,009,917
\end{tabular}
\end{center}
The number of monomials increases: fewer variables is not fewer expanded
terms. The full affine substitution map is rebuilt and checked.

For a polynomial over unrestricted integers, only the 3,960 semantic
inputs need explicit nonnegativity constraints. For each such input $x$,
append
\[
 x-a^2-b^2-c(c+1)=0.
\]
Every natural $x$ has this form: apply the three-square theorem to
$4x+1$, whose representation has one odd and two even coordinates, then
divide the identity by four after removing the odd square's constant.
Conversely $a^2+b^2+c(c+1)\geq0$ for integer $a,b,c$.
Internal circuit wires may remain unrestricted because their triangular
multiplication definitions force the unique nonnegative values after
semantic inputs are nonnegative. Adding the squares of these 3,960
quadratic residuals gives the literal integer-domain polynomial in
\path{fixed_base/equation_integer.json}:
\[
 \boxed{\text{27,916 unrestricted integer variables, exact degree four,
 758,275 monomials.}}
\]
Its maximum absolute coefficient is 320,307,372,009,918.
Both lack-of-solutions statements are provably equivalent in ZFC to
$\operatorname{Con}(\mathrm{ZFC})$. Neither count competes with the
low-variable branch; the advance is an explicit, independently auditable
ordinary equation from the direct proof source.

\subsection{Scope and limitations}

The direct equation is much larger than the low-variable Pell-based
constructions. It should not be described as a numerical record.
Its contribution is completion: no bounded quantifier, variable exponent,
factorial, binomial coefficient, unknown compiler output, or unspecified
coefficient remains in the final polynomial file. The fixed numeric proof
source, the general arithmetic representations, and the flattening theorem
are separate, auditable layers. Future source optimizations can now be
measured against an actual ordinary equation rather than an interface.


\section{The complete ten integer base formula}
The following definitions make the norm transformations interpretable without
confusing a circuit register with an extra variable. The ten unrestricted
unknowns are
\[
 t,g_0,c_0,k_0,l,w,x,y,z,m.
\]
The fixed H1868 encoding uses $\kappa=2\,134\,940$ and the 31 weights
$w_i$, 22 targets $q_j$, and three literal integer polynomials
$D_2,D_1,D_0$ in \path{norm/data/code.json} and
\path{norm/data/kernels.json}. Their actual coefficient bytes, including
all fixed source ports, are included and authenticated. Set
\begin{align*}
 b&=6t+4,&g&=bg_0,&C&=6c_0-1,&k&=b(3k_0+2),\\
 B_c&=2^\kappa b^2,&\ell_c&=\sum_iB_c^{w_i},&v_c&=\ell_c+g.
\end{align*}
\begin{align*}
 c_c&=D_2(B_c)v_c^2+D_1(B_c)v_c+D_0(B_c),&t_c&=\sum_jB_c^{q_j},\\
 N_0&=B_c^{117},&N_c&=B_c^{1986},\\
 \mu_c&=N_0-1-(b-1)\ell_c,&S_c&=g+N_0(2c_c+t_c),\\
 T_c&=\mu_c+N_0(B_c-2)t_c-2,&X&=N_c(S_c+T_c+1)-S_c-1,\\
 Y&=2^{34\,159\,017}bB_c^{2109},&\mathcal C&=bB_c^{116}.
\end{align*}
These powers have fixed integer exponents. The data files are complete fixed
coefficient lists, not unspecified parameters of an existence theorem.
Continue with the integer polynomial registers
\begin{align*}
 L&=lY,&U&=2LX,&V&=4gwY,&W&=bw,\\
 A&=U(V+1),&B_1&=2X+1,&S&=2A-5.
\end{align*}
\begin{align*}
 T&=3WC-2(W^2-1),&P_h&=(LX)UV,\\
 K&=X+1+k(2P_h-2),&Q&=(P_h^2-1)K^2+1,\\
 \Delta_g&=g(\mathcal C-g)K^2-\mathcal C^2(C-KL)^2,&\xi&=zx.
\end{align*}
The new guard is justified in the next section. With
$\operatorname{halfA}=(LX)(V+1)$, so $A=2\operatorname{halfA}$, put
\begin{align*}
 D&=(A^2-4)C^2+4,&E_0&=C^2\xi,&E&=DE_0,\\
 F&=4(A^2-4)C^4D\xi^2+1,&G&=(\operatorname{halfA}-1)(F-1)-1,\\
 H&=C+F(B_1-C+16Ey),&I&=(G^2-1)H^2+1,\\
 \rho&=T+Sy,&\Sigma&=\xi S^5,&
 \Theta&=(\rho^3-\rho S^2-S^3)(\rho^2+23S^2),\\
 r&=bQ,&a_0&=r+(9\Delta_g-3)(rz)^2,&a_1&=DFI,\\
 Z&=\Sigma m-\Theta.
\end{align*}
The base ordinary equation is
\[
 \bigl(Z^2-\Sigma^2(a_0+a_1)\bigr)^2-4\Sigma^4a_0a_1=0.
 \tag{N}
\]
Every displayed symbol other than the ten unknowns is a fixed polynomial
abbreviation. The exact optimized circuit, rather than this unoptimized
rendering of its formula, determines the operation count.

\subsection{Inherited semantic theorem and its use}
The retained baseline proof establishes the preceding source decoder,
the binomial bridge, signed-domain extraction and norm-saturated Lucas
certificate. Its key reliability statement is that, under the sign and
size conclusions forced by the guard, square $DFI$ recovers the required
Lucas index. The proof uses pairwise coprimality of $D,F,I$, Pell
classification, index separation and rank divisibility. Its completeness
statement supplies $\xi$ divisible by any prescribed positive integer and
permits every residue class of $y$ modulo $\xi$ while keeping $DFI$ square.
This is the full residue interpolation property, not merely existence of
one suitable $y$.

For arbitrary signed inputs, the two-root norm first gives square $a_0,a_1$
and $\Sigma\mid\Theta$. Coprime square-factor extraction and residues modulo
three establish that $b,Q$ are positive squares before guard signs are used.
The guard then supplies the approximation and label bounds needed by the
Lucas theorem and source decoder. Conversely canonical source witnesses,
a Pell solution for $a_0$, the stated interpolation theorem, and the
intersective quintic $(v^3-v-1)(v^2+23)$ produce an integer $m$ solving (N).
The full proof and earlier coefficient/compiler calibration are retained in
\path{norm/BASELINE_PROOF.md} and \path{inherited/round11/}; they are
mathematical dependencies, not assertions that imported formal files were
rebuilt. The new guard is substituted only after proving its signed
soundness and canonical completeness, as follows.


\section{A source-aware approximation guard}
Let $W=116$, $B=2^{2134940}b^2$, and $\mathcal C=bB^W$.
The canonical digit code is
\[
 g=\sum_{i=1}^{31}(p_i-1)B^{w_i},\qquad
 1\le p_i\le b,
\]
where the weights are distinct positive integers at most $W$.
The coordinate $a$ has weight $W$. Source equation~4 is
$a=E+(N+1)s$; all three coordinates on the right are positive, so
$a\ge3$. Consequently $g\ge2B^W$, and in particular $g\ge B^W$.
The usual digit bound gives $g<\mathcal C$: indeed
\[
 g\le(b-1)\sum_{j=1}^{W}B^j
 < (b-1)\frac{B^{W+1}}{B-1}<bB^W,
\]
since $B>b$. Moreover $B^W\ge b^2$, whence
\[
 g^3\ge B^{3W}\ge b^2B^{2W}=\mathcal C^2.
\]

Write $e=C-KL$ and replace the former approximation guard by
\[
 \boxed{\quad R=g(\mathcal C-g)K^2-\mathcal C^2e^2.\quad}
\]
For arbitrary integers with $\mathcal C>0$, the condition $R>0$
forces $K\ne0$ and $0<g<\mathcal C$. The identity
\[
 4g(\mathcal C-g)=\mathcal C^2-(\mathcal C-2g)^2
 \le\mathcal C^2
\]
then gives $4e^2<K^2$. Thus the new guard implies the previous
range and approximation conclusions, and in fact halves the permitted range.

Conversely take a canonical source code and the existing approximation
$|e|<|K|/(4g)$, with $K\ne0$. Since $\mathcal C-g$ is a positive integer,
\[
 \frac{R}{K^2}
 >g(\mathcal C-g)-\frac{\mathcal C^2}{16g^2}
 \ge g-\frac{g}{16}=\frac{15g}{16}>0.
\]
This proves completeness without a stronger approximation theorem or a new
coefficient bound. The argument depends only on the verified top digit and
existing radix and source conventions. Source encoding $g$ is not the
separate source coordinate also named $g$ (whose weight is $9$).


\section{Centering the final radical}
Here $B$ abbreviates the already defined $B_1$.
The inherited baseline equation uses ten unrestricted integer variables and has degree
$1\,245\,324$ and $480$ binary arithmetic operations.  The new short guard,
polarization of the final norm, exact scalar retiming, and reuse of
$uv=(LX)V$ in $\operatorname{halfA}=LX+uv$
and $P_h=Uuv$ give a
$471$-operation equation at the same degree.  A further idea trades a small
increase in circuit size for a lower degree: translate the final witness by
an integer polynomial approximating the largest radical.

\begin{center}
\begin{tabular}{lrr}
Literal circuit & Operations & Exact degree\\\hline
\path{integer10_polar.json}&471&1\,245\,324\\
\path{integer10_center0.json}&485&1\,179\,788\\
\path{integer10_center1.json}&496&1\,114\,252\\
\path{integer10_center2.json}&545&1\,080\,888
\end{tabular}
\end{center}
All four use ten unrestricted integers and the same actual fixed integer
leaves. These are cost--degree tradeoffs, not uniform improvements or claims
of minimality.

\subsection{A universal coordinate-change lemma}
Write $Z=\Sigma m-\Theta$, and let the previous final norm be
\[
 N=(Z^2-\Sigma^2(a_0+a_1))^2-4\Sigma^4a_0a_1.
\]
Suppose integer polynomials $P,J$, independent of $m$, satisfy $a_1=P^2+J$.
Define
\[
 T=Z+\Sigma P,\qquad
 U=Z^2+2\Sigma PZ+\Sigma^2(a_0-J),\qquad
 N_c=U^2-4\Sigma^2a_0T^2.
\]
Then $N_c(m)=N(m+P)$ identically. Indeed,
$U=T^2+\Sigma^2(a_0-a_1)$ and polarization gives the identity.
The substitutions $m\mapsto m\pm P$ are inverse bijections on integers,
including when $\Sigma=0$. Thus this change requires no new number theory,
nonzero assumption, or loss of signed solutions. The cancellation in $J$
is displayed explicitly instead of hidden in $a_1-P^2$.

\subsection{Two economical centers}
Keep the inherited registers
\[
 D=(A^2-4)C^2+4,\quad E_0=C^2\xi,\quad E=DE_0,
 \quad F=4(A^2-4)C^4D\xi^2+1,
\]
\[
 G=(A/2-1)(F-1)-1,\quad
 H=C+F(B-C+16Ey),\quad I=(G^2-1)H^2+1,
 \quad a_1=DFI.
\]
Here $A/2$ denotes the already available integer polynomial
$\operatorname{halfA}$, not a division operation in the circuit.
The first center takes
\[
 R_0=2AE,\qquad J_{DF}=D-16E^2.
\]
The second, stronger center takes
\[
 R_0=2AE_0(D-2C^2),\qquad
 J_{DF}=D+16E_0^2\bigl((3C^2-4)D-4C^2(C^2-1)\bigr).
\]
In either case $DF=R_0^2+J_{DF}$, hence
\[
 P=R_0GH,\qquad J=J_{DF}G^2H^2-DF(H^2-1)
\]
satisfy $a_1=P^2+J$. In the stronger version, $DF$ has degree six in
$A$, but its square defect has degree only two in $A$.

\subsection{A further correction with integral coefficients}
For the second center put $c_2=C^2$, $h=A/2$, and define
\begin{align*}
 L_*&=1+16E_0^2(3c_2-4),& F_0&=1+64E_0^2(c_2-1),\\
 \Omega_0&=-64E_0(c_2-1),& H_0&=C+F_0(B-C+\Omega_0y),\\
 V_0&=-F_0H_0,& J_f&=E_0(c_2A^2+4-8c_2),\\
 W_*&=B-C+16Ey.
\end{align*}
Use the following division-free definition:
\[
 Q_*=(2h^2-2)EH-A\bigl[J_f(C+(F+F_0)W_*)+4c_2F_0^2y\bigr].
\]
One verifies $4c_2\xi A Q_*=GH-V_0$. This is an identity, not a quotient
used to define $Q_*$. Its coefficients are integral in $h$, but need not
be integral in an unconstrained formal variable $A$.
Finally put
\begin{align*}
 P&=R_0GH+L_*Q_*,\\
 J_0&=-4(c_2-1)-256E_0^2(c_2-1)^2,\\
 J&=J_0(GH)^2+c_2L_*AGH\bigl(AV_0-4\xi(4-6c_2)Q_*\bigr)\\
 &\hspace{1cm}-(L_*Q_*)^2-DF(H^2-1).
\end{align*}
Then $a_1=P^2+J$ again. To prove it, write
$J_{DF}=c_2L_*A^2+J_0$ and
$R_0=2Ac_2\xi(A^2c_2+4-6c_2)$, expand, and use the identity for $Q_*$.
This establishes the third centered equation without excluding exceptional
values of $A,C,\xi$.

\subsection{Exact degree and verification boundary}
For the fixed source, $a=\deg A=16\,384$,
$\deg\Sigma=81\,922$, $\deg a_0=294\,382$, and
$\deg a_1=458\,818$. The first two centered degrees are
$72a+140$ and $68a+140$. The deepest center reaches
\[
 4\deg\Sigma+\deg a_0+\deg a_1=1\,080\,888.
\]
Here the cross-term $-4\Sigma^2a_0T^2$ dominates; further corrections to
this same radical alone do not remove that obstruction.
The independently authored checker verifies the exact polynomial
identities, integrality in $\operatorname{halfA}$, actual numeral payloads,
full source-to-equation replays, topology, liveness, and operation counts.
Nonzero highest-degree coefficients modulo three primes certify attainment
of the stated degrees. Finite replay is an audit of transcription, not a
replacement for the universal identity or the inherited consistency
calibration. These results are mathematical proofs, not new Lean/Rocq
certificates.


\section{Independent audit of the centered integer equations}
An independently written checker imports no producer builder, circuit helper,
evaluator, or checker. It validates all actual fixed numeral payloads and
source-port bindings, reconstructs all $22$ source target identities across
$20\,466$ kernel interactions, and checks every named literal register by
exact polynomial coefficient cancellation. It verifies the universal square
completions and norm translation over integer polynomial rings. In particular,
the deepest correction is certified in
$\mathbb Z[\operatorname{halfA},C,\xi,B,y]$ after the substitution
$A=2\operatorname{halfA}$, not admitted as a divided expression in an
unrestricted independent $A$.

The four independent operation counts are $471,485,496,545$, with respective
exact total degrees $1\,245\,324$, $1\,179\,788$, $1\,114\,252$, and
$1\,080\,888$. The final exact reuse of $\operatorname{halfU}V$ in the $\operatorname{halfA}$ and $\operatorname{Phalf}$ registers saves one operation without changing their polynomials. Each exact-degree claim has three nonzero leading-coefficient
certificates using the actual fixed integers. The checker also performs
$448$ signed and degenerate modular replays per candidate, including moduli
$2$ and $3$. Normal and optimized Python execution use explicit exceptions,
not assertions. These finite checks are regression tests, not substitutes for
universal mathematical proofs.

Two logical distinctions matter. First, the square-factor and residue
argument recovers positive squares $b,Q$, and thus positive
$\mathcal C=bB^{116}$, before the short guard is used to infer
$0<g<\mathcal C$ and $4(C-KL)^2<K^2$. There is no sign assumption hidden in the
ten unrestricted integer inputs. Second, replacing the guard preserves
existence of solutions through all-signed soundness and canonical-witness
completeness; it does not assert a bijection between old and new arbitrary
witness tuples. Only the subsequent triangular changes
$m\mapsto m+P$ are literal integer bijections, with inverse $m\mapsto m-P$,
even when $\Sigma=0$.

The inherited source/compiler/refutation calibration and Pell--Lucas
reliability and interpolation theorems remain explicit ordinary-mathematical
dependencies. No new formal-kernel certificate, historical priority, or
optimality is claimed. Full receipts, circuit hashes, exact identities, and
a proof narrative appear in the \path{norm_audit} component.


\section{Eleven positive variables using two small gaps}
\label{sec:positive211}
The new polynomial has eleven positive variables, degree $211$, and $113$
live binary arithmetic operations.  It improves the incoming centered
$257/114$ construction in both degree and evaluation cost.  The earlier
$277/111$ point still uses two fewer operations.  Centering itself is not
new: our contribution is a second center, a cheaper gap bound, and a
cofinality argument that permits a useful change of the auxiliary witness.

\subsection{The literal polynomial}
Use the complete compact44 coefficient tuple
$\beta,K_0,M_C,M_F$, with $M_F$ already shifted, as bound in
\path{positive/native/ACTIVE_BINDING.json}.  The unknowns are
\[
 j,F,\alpha,i,T,X,Y,n,\eta,\zeta,\gamma\in\mathbb Z_{>0}.
\]
Set
\begin{align*}
 J&=3j+1,&q&=\beta J+1,&C&=q-F-\alpha,\\
 E&=XY,&k&=\eta+\zeta,&c&=kY+\eta,\\
 a&=E+Y,&H&=4a+3,&\Delta&=a^2+H=(a+2)^2-1,\\
 \ell&=EkY,&I&=ic^2,&D_f&=1+\Delta I^2,\\
 Q&=\Delta(D_f-1),&\delta&=Q-1,&D_\tau&=\ell(\ell+k)+1.
\end{align*}
The coefficient and transport registers are
\begin{align*}
 R&=((q-1)(q-F)+\alpha)(q^2-1)+(M_C+qM_F)J,\\
 M&=q^3(q-1),\qquad B_t=(qK_0+X)C-qF.
\end{align*}
Here are the new centers and the small transport gap:
\begin{align*}
 T_o&=T+q+\alpha,& b&=cT_o-R,&V&=D_f b-c,\\
 D&=(a+2)I,&L&=kT_o+F,&W&=1+(q-1)L,\\
 \overline Z&=M(\ell+k+n+L+V+D)+(q-1)Y-q^3WB_t.
\end{align*}
Keep the adaptive main carrier
\begin{align*}
 r_0&=k-R-1,&S&=4qY,&U&=qS,\\
 D_*&=U(X+ac)+H(U\gamma+SX-r_0),\qquad
 A_1=D_*^2-U^2(1+\Delta c^2).
\end{align*}
Define the three-root norm without any division:
\begin{align*}
 A&=M^2D_\tau,\quad B=M^2D_f,\quad C_N=M^2(QV^2-1),\\
 T_N&=\delta^2AB,\quad R_N=\delta C_N,\\
 J_N&=\delta(\overline Z^2-A-B)-C_N,\\
 N_c&=J_N^2-4[T_N+R_N(A+B)],\\
 N_8&=N_c^2-64T_NR_N\overline Z^2,\\
 \boxed{P_{11}}&\boxed{=N_8+A_1+N_8A_1=0.}
\end{align*}
All registers are abbreviations, not additional quantified coordinates.
\par The exact circuit is \path{positive/p11_degree211.json}.

\subsection{Recover divisibility before using positivity}
Write, only in the proof,
\[
 D_y=\frac{QV^2-1}{Q-1},\qquad
 \tau=\sqrt{D_\tau},\quad f=\sqrt{D_f},\quad y=\sqrt{D_y}.
\]
The inherited unconditional separation lemma applies for every integral
$b$, without assuming $C>0$ or $V>0$.  Indeed $I\ge c^2$, $c\ge3$ imply
$1<D_\tau<D_f<Q$ and $|V|>1$.  If $b=0$, then $V=-c$ and
$f>\tau+y$; otherwise $y>f+\tau$.  The radicands are distinct, since
$D_y=D\in\{D_\tau,D_f\}$ would imply $Q(D-V^2)=D-1$, with
$0<D-1<Q$.  Thus no signed sum vanishes.

The norm equation gives
$\overline Z/M=\tau_*+f_*+y_*$ with the appropriate signed roots.
These roots are rational.  An explicit extraction identity is as follows:
for squares $a_1,a_2,a_3$ and nonzero sum $s=x_1+x_2+x_3$, put
\[
 u=\frac{s^2-a_1-a_2-a_3}{2},\qquad
 v=\frac{u^2-a_1a_2-a_1a_3-a_2a_3}{2s}.
\]
Then $x_i=a_i(a_i+u)/(a_is+v)$; the denominator equals
$x_i(x_i+x_j)(x_i+x_k)$ and cannot vanish by distinctness.
The first two roots are integral because their squares are integers.
The adjacent-Pell rational-integrality lemma makes $y_*$ integral from
$QV^2-(Q-1)y_*^2=1$.  For completeness, write $y_*=a/m$ reduced;
$m^2\mid Q-1$, and the norm-preserving descent
\[
 V'=(2Q-1)V-2(Q-1)y_*,\qquad
 y'=(2Q-1)y_*-2QV
\]
has integral $V'$, unchanged reduced denominator $m$, and
$0<|V'|<|V|$ when $|V|>1$.  The terminal case $|V|=1$ forces $m=1$.

Therefore $M\mid\overline Z$, and reduction modulo $q^3$ gives
$q^3\mid Y$.  This is the key dependency order: no transport-gap sign
or positivity of $C$ has been used yet.

\subsection{A degree-two bound for the auxiliary gap}
The authentic masks give $0<(M_C+qM_F)J<2q^2$ and $q\ge16$.  Directly,
\[
 -Fq^3<R<q^3(q+\alpha).
\]
For the upper inequality, subtract the bound
$q(q-1)(q^2-1)+\alpha(q^2-1)+2q^2$ from $q^3(q+\alpha)$;
the difference is $q^3-q^2-q+\alpha(q^3-q^2+1)>0$.
Since $c>Y\ge q^3$, we get
\[
 c(q+\alpha)>R,\qquad b>cT\ge c,\qquad V>1.
\]
Also $\delta>(\Delta-1)D_f$, $c<2kY$, and
$\Delta-1=(a+2)^2-2>4Y^2$.  Consequently
\begin{align*}
 0<y-V
 &=\frac{V^2-1}{\delta(y+V)}
 <\frac{b}{\Delta-1}\\
 &<\frac{cT_o+Fq^3}{\Delta-1}<kT_o+F=L.
\end{align*}
The second center satisfies
\[
 D^2-D_f=I^2-1>0,\qquad 0<D-f=\frac{I^2-1}{D+f}<I.
\]
Thus the huge root $y$ is transported only through $y-V$, and the strong
root $f$ through the favorable gap $D-f$.

\subsection{Soundness and cofinal completeness}
Recover positive witnesses
\[
 u=y-V,\qquad
 r=n+L+V+D-f_*-y_*,\qquad f_{\rm old}=|f_*|,\quad i_{\rm old}=i.
\]
Indeed $r\ge n+L-(y-V)+D-f>0$.  The weighted first carrier with
this $r$ and weight $L$ is $M\tau_*$.  The strong and auxiliary residuals,
and $A_1$, vanish.  The inherited weighted-source theorem therefore yields
a finite ZFC refutation.  The signed roots, not their absolute values,
must be used inside $r$.

Conversely, canonical source witnesses give the outer coordinates and a
normalized strong pair $f,I$, and the inherited auxiliary lemma supplies
one positive pair $V_0,y_0$ with
\[
 QV_0^2-(Q-1)y_0^2=1,\qquad V_0=f^2(cT_0-R)-c.
\]
The matrix
\[
 B_Q=\begin{pmatrix}2Q-1&2(Q-1)\\2Q&2Q-1\end{pmatrix}
\]
has determinant one, positive entries, and preserves the norm.
Some positive power is the identity modulo $h=cf^2$, by finiteness of
$\operatorname{SL}_2(\mathbb Z/h\mathbb Z)$.  Its iterates give unbounded
positive pairs with $V_j\equiv V_0\pmod h$, so
\[
 T_j=(V_j+c+Rf^2)/(cf^2)
\]
is integral and unbounded.  Choose $T_o=T_j>q+\alpha$ with
$L=kT_o+F>I$.  Then $T=T_o-q-\alpha>0$.
The canonical weighted carrier supplies
$r>qW-2k>2L>L+I$; hence
\[
 n=r-L-V-D+f+y=r-L-(D-f)+(y-V)>0.
\]
It has $\overline Z=M(\tau+f+y)$, and the unchanged adaptive carrier
supplies positive $\gamma$.  This proves completeness without assuming
an unjustified size condition on the original auxiliary witness.
Finally, $N_8$ and $A_1$ are squares modulo four, so their unit-product
finalizer vanishes exactly when both vanish.  Negating solvability gives
$\operatorname{Con}(\mathrm{ZFC})$ in ZFC, with the same explicitly
inherited compiler and calibration boundary.

\subsection{Exact degree rather than a generic gate bound}
Put $z=\overline Z-M(V+D)$, $e=I^2-1$, and
$H_0=z^2+2MDz+M^2e-A$.  Then
\[
 J_N=\delta[H_0+2MV(z+MD)]-M^2(V^2-1).
\]
The identity $B=M^2(D^2-e)$ cancels the apparent degree-102 part of
$N_c$.  Its degree-99 head is $8Q_0^2M_0^3V_0^2D_0z_0$, and the
norm has degree-198 head
\[
 64Q_0^4M_0^6V_0^4D_0^2(z_0^2-M_0^2D_{\tau,0}).
\]
With $q_0=3\beta j$, $k_0=\eta+\zeta$,
$L_0=k_0(T+q_0+\alpha)$ and
$B_{t,0}=(K_0q_0+X)(q_0-F-\alpha)-q_0F$, we have
\[
 z_0=M_0(\ell_0-L_0B_{t,0}),\qquad
 z_0^2-M_0^2D_{\tau,0}
 =M_0^2L_0B_{t,0}(L_0B_{t,0}-2\ell_0)\ne0.
\]
Multiplying by the degree-13 adaptive-carrier head shows that the
coefficient of
\[
 j^{48}i^{18}\eta^{39}X^{30}Y^{67}F^4T^2\alpha^2\gamma
\]
is $8192(3\beta)^{48}>0$.  This certifies exact degree 211 on the
actual coefficient slice, independently of $K_0,M_C,M_F$.

\subsection{A materialized incoming degree 70 companion}
The manuscript-only incoming \path{report.zip} already proposes a
12-positive-variable equation of degree 70 with 100 operations.  Its advertised
JSON was absent.  We supply it as
\path{positive/p12_degree70.json}, with independent checks matching
all metrics, but do not claim its theorem as new here.  It replaces the
normalized strong norm by the ordinary strong norm after an early index
bound licenses ordinary rank recovery.  Its exact residual degrees are
$15,13,42$.  The full proof is retained verbatim in
\path{positive/interfaces/incoming_report.tex}; the crucial endpoint
is the ordinary source, not a possibly nonintegral division of the strong
multiplier by $\Delta$.


\section{Low-degree integer equations by delayed transport elimination}
The new low-degree frontier point established here is
\[
 (n,d)=(21,16)\qquad\text{over unrestricted integers}.
\]
The completed additional charts $(20,20),(19,32),(18,88)$ improve the
initially inspected incoming family, but not the later Library supplement,
which already supplies $(20,20),(19,28),(18,32),(17,88)$.
Thus the two $(20,20)$ charts are alternative derivations, and the
$(19,32),(18,88)$ charts are dominated comparisons, not new frontier records.
They are retained to expose the reversible substitutions. These are
variable/degree results, not operation-count or global-minimality claims.
Each equation is the sum of squared literal residuals in its packet.

\subsection{An exact pivot that becomes inexpensive later}
Use the complete source34 tuple, including the already shifted $M_F$,
and the fixed scales $A=2^{1113}$ and $B=2^{666104}$. Set
\[
 q=\beta(3J_3+1)+1,\quad W=4(3w_3+1),\quad
 C=q-F-\alpha.
\]
The parent transport row is
\[
 (K_0+W)C-F-z(q-1)=0.
\]
After replacing $\alpha$ by $C$, it is monic in $F$. Hence
\[
 F=(K_0+W)C-z(q-1),\qquad \alpha=q-F-C
\]
are integral polynomial restorations; neither $q-1$ nor any fixed
coefficient is divided out. The determinant guard
\[
 T=AFBC-(q-AF-BC)^2
\]
then has degree four, and the packing row degree five. These costs are
small compared with the larger degree budgets in the later projection
families. This is why delaying the pivot saves a variable without
raising their final degree.

\subsection{The primary point with 21 integers at degree 16}
Start from the exact 24-variable joint-positivity parent. Write
$\Delta=a^2+4a+3$ and $S=I^2-\Delta(f^2-1)$. Besides the transport
pivot, eliminate $V=fo-c$, shear the auxiliary root
$y_{\rm old}=fo+y$, and replace its old row by
\[
 A_{\rm new}=(fo+y)^2-
 \Delta(f^2-1)(c+y)(2fo+y-c)-1.
\]
The identity
\[
 A_{\rm new}=A_{\rm old}
       +S\bigl(y_{\rm old}^{\,2}-V^2\bigr)
\]
is reversible while retaining $S=0$. Now eliminate
$I=(2r+1)c^2$. The three deleted monic definitions leave
$24-3=21$ unknowns and $12-3=9$ rows. The strong row has degree six,
the centered auxiliary row seven, and the joint guard eight.
The coefficient certificate is
\[
 [w_3^8C^8]P_{21}=12^8A^8>0.
\]
Thus the degree is exactly sixteen on the actual fixed coefficient
slice. All integer signs and all parent witnesses are covered.

\subsection{The smaller-variable tradeoffs}
For these endpoints retain the original auxiliary row instead. Together
with the transport pivot and the monic $V,I$ eliminations, apply
\[
\begin{array}{c|l|c|c}
 n&\text{additional restorations}&\text{rows}&d\\\hline
20&a=E+Y&8&20\\
19&E=XY,\ a=(X+1)Y&7&32\\
18&E=XY,\ a=(X+1)Y,\ \eta=c-Yk&6&88
\end{array}
\]
where $X=qW$, $Y=(2v+1)q^3$, and $k=1-V+jc+hE$.
These are ordinary polynomial substitutions, not divisibility
relaxations. The selected exact leading coefficients are
\begin{align*}
[f^4v^4J_3^{12}]P_{20}&=2^4 3^{12}\beta^{12},\\
[\eta^4h^4w_3^4v^4J_3^{16}]P_{19}&=2^{12}3^{20}\beta^{16},\\
[h^8w_3^8v^{16}J_3^{56}]P_{18}&=2^{32}3^{64}\beta^{56}.
\end{align*}
They are strictly positive for the materialized tuple. The eighteen
unknowns are exactly
\[
 J_3,z,f,h,r,v,w_3,\tau,y,\gamma,o,j,c,t,p,q_{\rm pos},r_{\rm pos},C.
\]
A different twenty-variable chart uses one free scale $Y_0$ with
$E=XY_0$, $a=(X+1)Y_0$, retaining
$Y_0-(2v+1)q^3=0$. It again has degree twenty, with selected coefficient
$36^4\beta^4$ at $\eta^4h^4J_3^4w_3^4Y_0^4$.

\subsection{Correctness and checks}
Projection and the displayed polynomial restorations are inverse on
integer zero sets; the auxiliary shear and row operation are explicitly
reversible. Composing them with the pinned parent's mathematical
native/Pell compiler theorem yields, in ZFC, nonexistence of an integer
zero if and only if $\operatorname{Con}(\mathrm{ZFC})$. This is not a
new kernel-checked certificate or a proof of consistency.

The package includes all six actual fixed coefficients, complete
literal residuals and ordered variable lists, the inherited proof
chain, and a separate full proof. A newly authored builder verifies
homogeneous-top coefficient certificates. An independent standard-library
replay checks all restored rows on 750 signed modular assignments and
authenticates all coefficient payloads, in normal and optimized Python.
Finite replay is a transcription check, not the unbounded equivalence
proof. No downloaded generator is executed.


\section{Exact equations and reproducible evidence}
The release contains three explicitly defined ordinary-polynomial formats.
They differ only in storage, not in logical status.
\begin{enumerate}
\item A sparse list entry $[c,[i_1,\ldots,i_s]]$ denotes
$cX_{i_1}\cdots X_{i_s}$. Repeated indices repeat factors; the empty list
denotes the constant term. Summing the complete list in
\path{direct/fixed_base/equation_integer.json} defines the fully collected
quartic, with zero-based indices from 0 to 27,915.
\item A residual packet defines precisely $\sum_jr_j^2$, where every
$r_j$ and every intermediate abbreviation is explicitly listed and all
exponents are fixed. The packet's ordered variable list is exhaustive.
This is used by the 21-integer equation and its alternative charts.
\item A literal arithmetic DAG defines the polynomial at its designated
output. Only fixed integer leaves and the declared unknowns enter; every
other node is addition, subtraction or multiplication of earlier nodes.
All fixed numeral payloads are present. It is a compact exact polynomial,
not an existential statement that a suitable polynomial could be generated.
This is used by the ten-integer and positive-variable equations.
\end{enumerate}
The direct sparse file itself is included, not merely a generator or a hash.
To keep the archive small, redundant intermediate files are
regenerated before auditing. Their SHA-256 values are fixed in the replay
record. The natural equation is already completely specified by its literal
quadratic residual list.

\subsection{Verification contract}
The validated runtime is Debian 13 Linux x86\_64, Python 3.12.14,
SymPy 1.14.0 and zlib 1.3.2. Exact compressed-byte regeneration is claimed
for this environment, not all Python releases; another gzip implementation
may encode the same integer bytes differently. The complete runtime record
is \path{RUNTIME.json}.

The top-level \path{verify.py} authenticates the manifest and copies the
archive into an isolated temporary directory before executing the explicit
\path{REPLAY.json} allowlist. Only newly authored builders, checkers and
retained authored packing checks are run. Static upstream Python, Lean,
machine data and proof references remain inert. Ordinary and optimized
Python runs are both required; the checks use explicit exceptions and do
not depend on enabled assertions. Rebuilding must reproduce every literal
equation and actual coefficient payload byte for byte.

Independent exact checks cover all Horn branches and constructor wires,
every emitted Pell/power/division certificate, the binomial/factorial
compositions, the local-source and bounded-CRT bindings, every affine
substitution, and the complete collected direct quartic. The other audits
check polynomial square completions over integer rings, source coefficient
bindings, operation counts, live nodes, exact leading coefficients, signed
restorations, and the positive-domain cofinality proof. Mutation negative
controls verify that selected damaged circuits, source rows, fixed payloads
and degree claims are rejected. The clean extracted ZIP is replayed as a
separate packaging test.

A finite modular regression checks a transcription, not a universal theorem.
An exact nonzero leading coefficient proves attainment of a degree once its
symbolic upper bound has been established. An exact identity check proves
that identity but not the imported proof-calculus calibration. These distinct
claims are recorded separately in the component audit reports.

\subsection{The scope of the consistency claim}
For the direct construction, the calibration is the fixed numerical theorem
\path{KEndpoint.c222_iff_zf_refutation} followed by
\path{ConstructibleChoice.zf_refutation_iff_zfc_refutation}. The numerical
constructors and rules are retained in the pinned source files. Their
mathematical proof uses ordinary syntactic proof transformations and the
ZF/ZFC equiconsistency bridge. It does not simulate K215.

The small-variable branches retain their earlier native coefficient/compiler,
Pell and machine-based calibration. This report does not pretend that a new
coordinate transformation removes those dependencies. Both routes yield
ordinary mathematical biconditionals formalizable in ZFC; neither supplies
an independently compiled end-to-end object-language ZFC derivation or a new
Lean/Rocq build. No implication concerning the truth of consistency is
inferred from a search that finds no zero.

\section{Further research with concrete stopping points}
\begin{enumerate}
\item The direct quartic's main cost is generic exponentiation certification.
A specialized proof of the fixed-base or factorial subgraphs should be judged
by a regenerated ordinary equation, not by the size of an uncompiled source.
The current call ledger gives a reproducible starting point.
\item The common witness bound must remain independent of the label bound.
Identifying them would make a CRT bound depend circularly on the accumulator
it bounds. A successful replacement must prove a noncircular completeness
construction for arbitrary finite proofs.
\item Centering only the final large radical reaches a cross-term obstruction
at degree 1,080,888 in the present ten-variable coordinates. Further degree
reduction requires changing that cross term or another source register, not
merely extending the same formal square-root expansion.
\item Combining the Library supplement's CRT transport chart with the
centered auxiliary row gives another 21/16 chart, but no better completed
metric. Eliminating its retained scale raises the relevant degree; removing
its rank register exposes a larger square. An improvement must overcome
those explicit degree obstructions while preserving an integer inverse.
\item The 11-positive construction depends on selecting an auxiliary Pell
solution in an unbounded congruence orbit. Any sharper gap estimate must
preserve that cofinality and the order ``integrality, divisibility, then
positivity.'' Assuming a favorable sign before extracting the roots would
invalidate soundness.
\end{enumerate}
These are research questions rather than claims of additional completed
results. The present deliverables stop at the exact equations, proofs and
replayable checks described above.

\section*{Sources and attribution}
\addcontentsline{toc}{section}{Sources and attribution}
The principal source is Vladimir Reshetnikov's
\href{https://github.com/VladimirReshetnikov/ProveIt}{ProveIt repository},
with the current comparison pin and per-file authenticated source records
in the archive. Earlier local proof packets are retained with their original
provenance, licenses and notices. The arithmetic representation sources
include the finite Matiyasevich Pell criterion, the explicit power graph,
factorial/binomial identities and the bounded-CRT certificate. The inherited
Lucas classification, index separation and interpolation interfaces are
attributed in the retained proofs to Zhi-Wei Sun. K215/Wade licensing notices
apply only to the small-variable source branch and are preserved.

This is AI-assisted research prepared by dot, an OpenAI-powered assistant,
for Vladimir Reshetnikov. New transformations and checks are not presented
as peer-reviewed results, human-authored work or endorsements by upstream
contributors. No claim of historical priority or global minimality is made.
\end{document}
